%%%%%%%%%%%%
%
% $Autor: Wings $
% $Datum: 2019-03-05 08:03:15Z $
% $Pfad: TemplateSensor $
% $Version: 4250 $
% !TeX spellcheck = en_GB/de_DE
% !TeX encoding = utf8
% !TeX root = filename 
% !TeX TXS-program:bibliography = txs:///biber
%
%%%%%%%%%%%%

% Structure
\chapter{Konfiguration PlatformIO für Arduino MKR WAN 1310}
Um die Kommunikation mit dem Arduino MKR1310 herstellen zu können, muss die Datei \FILE{platformio.ini} die korrekten Konfigurationseinstellungen enthalten. Diese lassen sich bereits bei der Erstellung eines neuen Projekts vorkonfigurieren, indem das passende Entwicklungsboard ausgewählt wird.

\vspace{0.5cm}
Um dies sicherzustellen, sind die folgenden Schritte auszuführen:

\begin{enumerate}
	\item Neues Projekt erstellen:
	\begin{center}
		\includegraphics[width=0.7\linewidth]{PlatformIO/PIONewProject}
		\captionof{figure}{Neues Projekt erstellen.}
		\label{fig:PIONewProject}
	\end{center}
	\item Den Projektnamen, den Boardtyp und den Speicherort angeben.
	\begin{center}
		\includegraphics[width=0.7\linewidth]{PlatformIO/PIOProjectDet}
		\captionof{figure}{Den Namen, den Boardtyp und den Speicherort angeben.}
		\label{fig:PIOProjectDetails}
	\end{center}
	
	Als „Board“ ist „Arduino MKR WAN 1310“ auszuwählen, da dieses aktuell verwendet wird. Die Datei \PYTHON{platformio.ini} wird anschlie{\ss}end automatisch für dieses Board vorkonfiguriert.
	
	
	Hier kann auch der Speicherort festgelegt werden. Standardmä{\ss}ig werden alle Projekte unter \texttt{/Users/../Documents/PlatformIO/Projects} gespeichert.
	
\end{enumerate}
Somit ist die Konfiguration abgeschlossen. Weitere Parameter zur Konfiguration und zu den Bibliotheken können zu einem späteren Zeitpunkt in der Datei \PYTHON{platformio.ini} angepasst werden.